\chapter{Los alcoholes: activar el grupo OH}\label{ch:b1:alcohol-activation}

Los alcoholes están en todas partes: se fabrican por toneladas a partir
de alquenos y por fermentación, y los da cada adición de Grignard del
capítulo anterior. Sin embargo, un alcohol no sufre una sustitución como
un halogenoalcano: el grupo OH tendría que salir como ion hidróxido, una
\omterm{def:b1:predominant-reaction:strong-acid}{base fuerte} y un \omterm{def:b1:electronic-effects:nucleophile}{grupo saliente} muy malo
(\cref{ch:b1:nucleophilic-substitution}). Los químicos tienen tres
respuestas. Si se arranca el protón del OH, el \omterm{def:b1:alcohol-activation:alkoxide}{alcóxido} se convierte en
un \omterm{def:b1:electronic-effects:nucleophile}{nucleófilo} potente. Si se añade un protón, el \omterm{def:b1:electronic-effects:nucleophile}{grupo saliente} pasa a
ser el agua. O bien se sustituye el hidrógeno por un grupo que convierte
el oxígeno en parte de un \omterm{def:b1:electronic-effects:nucleophile}{grupo saliente} excelente, sin tocar el carbono.
Este capítulo sigue las tres vías, y las elecciones de mecanismo y de
estereoquímica que las acompañan.

\begin{recall}
La SN1, la SN2 y su estereoquímica (\cref{ch:b1:nucleophilic-substitution});
la E1, la E2 y la \omterm{prop:b1:elimination:zaitsev}{regla de Zaitsev} (\cref{ch:b1:elimination}); las
\omterm{def:b1:predominant-reaction:acidity-constant}{constantes de acidez} y el método de la \omterm{def:b1:predominant-reaction:predominant-reaction}{reacción predominante}
(\cref{ch:b1:predominant-reaction}); el efecto de los grupos alquilo sobre
la acidez de los alcoholes
(\cref{exo:b1:electronic-effects:11}).
\end{recall}

\section{Los alcoholes como ácidos: los alcóxidos}

\begin{definition}[Alcóxido]\label{def:b1:alcohol-activation:alkoxide}
Un \emph{alcóxido}\index{alcoxido@alcóxido} es la base conjugada \ce{RO-} de un alcohol
\ce{ROH}: metóxido \ce{CH3O-}, etóxido \ce{CH3CH2O-}, terc-butóxido
\ce{(CH3)3CO-}. Los alcóxidos son \omterm{def:b1:predominant-reaction:strong-acid}{bases fuertes} y buenos \omterm{def:b1:electronic-effects:nucleophile}{nucleófilos}.
\end{definition}

\begin{proposition}[Acidez de los alcoholes]\label{prop:b1:alcohol-activation:alcohol-acidity}
Los alcoholes son ácidos más débiles que el agua o casi tan débiles como
ella: $\mathrm{p}K_a$ 15.5 (metanol), 15.9 (etanol), 17.1 (propan-2-ol),
19.2 (2-metilpropan-2-ol), frente a 14.0 para el par \ce{H2O}/\ce{OH-}.
Por eso el hidróxido solo los desprotona en parte; una conversión
completa en el \omterm{def:b1:alcohol-activation:alkoxide}{alcóxido} necesita una base más fuerte o un metal alcalino.
\end{proposition}
% ledger: pka:methanol, pka:ethanol, pka:2-propanol, pka:tBuOH

\begin{proof}
La reacción \ce{ROH + OH- <=> RO- + H2O} tiene $K = 10^{14.0 -
\mathrm{p}K_a(\ce{ROH})}$, entre $10^{-1.5}$ y $10^{-5}$: desfavorable. El
sodio metálico reduce el protón, \ce{2ROH + 2Na -> 2RO- + 2Na+ + H2}, y el
ion hidruro del hidruro de sodio lo toma de forma irreversible,
\ce{ROH + H- -> RO- + H2}: el hidrógeno se escapa, y la reacción es
completa sea cual sea su constante de equilibrio.
\end{proof}

\begin{inthelab}[El hidruro de sodio]
El hidruro de sodio se vende como un polvo gris disperso en aceite
mineral, que lo protege de la humedad del aire. Reacciona violentamente
con el agua, desprendiendo hidrógeno inflamable; se pesa deprisa, se
añade en porciones al alcohol seco bajo gas inerte, y el burbujeo de
hidrógeno muestra el avance de la desprotonación. Al final, el hidruro en
exceso se destruye añadiendo lentamente un alcohol, nunca agua.
\end{inthelab}

\section{La síntesis de Williamson de éteres}

\begin{definition}[Síntesis de Williamson]\label{def:b1:alcohol-activation:williamson}
La \emph{síntesis de Williamson}\index{sintesis de Williamson@síntesis de Williamson} prepara un éter
$\mathrm{R{-}O{-}R'}$ mediante una reacción SN2 de un \omterm{def:b1:alcohol-activation:alkoxide}{alcóxido} \ce{RO-}
sobre un halogenoalcano (o un \omterm{def:b1:alcohol-activation:sulfonate}{éster sulfonato}) $\mathrm{R'{-}X}$.
\end{definition}

\begin{omfigure}
\setchemfig{atom sep=2.1em}
\schemestart
\chemfig{CH_3CH_2-@{o}O^{-}}\,\chemfig{Na^{+}}
\+
\chemfig{@{c}CH_3-[@{ci}]@{i}I}
\arrow{->}
\chemfig{CH_3CH_2-O-CH_3}
\+
\chemfig{Na^{+}}\,\chemfig{I^{-}}
\schemestop

\medskip
\chemmove{\draw[->, thick, shorten <=2pt, shorten >=2pt](o) .. controls +(80:8mm) and +(100:8mm) .. (c);
          \draw[->, thick, shorten <=2pt, shorten >=2pt](ci) .. controls +(90:5mm) and +(90:5mm) .. (i);}

{\small \omterm{def:b1:alcohol-activation:williamson}{Síntesis de Williamson} del metoxietano: el etóxido ataca el
carbono del yodometano por el lado opuesto al yodo (SN2).}
\end{omfigure}

\begin{method}[Elegir los dos participantes]\label{met:b1:alcohol-activation:williamson-choice}
Un éter $\mathrm{R{-}O{-}R'}$ puede cortarse a uno u otro lado del
oxígeno.
\begin{enumerate}
  \item Se escriben las dos parejas posibles: \ce{RO-} con
        $\mathrm{R'X}$, y $\mathrm{R'O^-}$ con \ce{RX}.
  \item Se conserva la pareja en la que el haluro es metílico o primario:
        la SN2 necesita un carbono no impedido.
  \item Se descarta una pareja con un haluro terciario: el \omterm{def:b1:alcohol-activation:alkoxide}{alcóxido}, una
        \omterm{def:b1:predominant-reaction:strong-acid}{base fuerte}, eliminaría (E2) en su lugar.
  \item Con un haluro secundario, cabe esperar una mezcla de éter y
        alqueno.
\end{enumerate}
\end{method}

\begin{example}[Un éter asimétrico]\label{ex:b1:alcohol-activation:williamson}
El 2-metoxi-2-metilpropano (metil terc-butil éter): el terc-butóxido de
sodio con yodometano lo da limpiamente; el metóxido de sodio con el
bromuro terciario, el 2-bromo-2-metil\-propano, da 2-metilpropeno por E2.
Voluminoso del lado del \omterm{def:b1:alcohol-activation:alkoxide}{alcóxido}, pequeño del lado del haluro.
\end{example}

\section{Activar el grupo OH}

\begin{definition}[Ésteres sulfonato]\label{def:b1:alcohol-activation:sulfonate}
Un \emph{éster sulfonato}\index{ester sulfonato@éster sulfonato} $\mathrm{R{-}O{-}SO_2{-}R'}$ se forma a
partir de un alcohol y un cloruro de sulfonilo $\mathrm{R'SO_2Cl}$ en
presencia de una base (piridina). El \emph{tosilato}\index{tosilato},
\ce{R-OTs}, procede del cloruro de 4-metilbencenosulfonilo (cloruro de
tosilo), y el \emph{mesilato}\index{mesilato}, \ce{R-OMs}, del cloruro de
metanosulfonilo. El anión sulfonato $\mathrm{R'SO_3^-}$, base conjugada de
un \omterm{def:b1:predominant-reaction:strong-acid}{ácido fuerte} y estabilizado por resonancia sobre tres oxígenos, es un
\omterm{def:b1:electronic-effects:nucleophile}{grupo saliente} excelente.
\end{definition}

\begin{omfigure}
\setchemfig{atom sep=2em}
\schemestart
\chemfig{R-O-H}
\+
\chemfig{Cl-S(=[2]O)(=[6]O)-[:0]*6(-=-(-CH_3)=-=)}
\arrow{->[piridina]}
\chemfig{R-O-S(=[2]O)(=[6]O)-[:0]*6(-=-(-CH_3)=-=)}
\schemestop
\par\vspace{6pt}

{\small Tosilación de un alcohol. El enlace O--H se sustituye por O--S;
el enlace C--O del alcohol no se toca. La piridina capta el HCl formado.}
\end{omfigure}

\begin{proposition}[La configuración en la tosilación y la sustitución]\label{prop:b1:alcohol-activation:tosylation-retention}
La tosilación de un alcohol sobre un carbono estereogénico conserva su
\omterm{def:b1:stereochemistry-in-depth:isomers}{configuración} (retención); una sustitución SN2 posterior del \omterm{def:b1:alcohol-activation:sulfonate}{tosilato} la
invierte. Las dos etapas juntas sustituyen el OH por un \omterm{def:b1:electronic-effects:nucleophile}{nucleófilo} con
inversión.
\end{proposition}

\begin{proof}
En la tosilación, los enlaces que cambian son O--H y S--Cl: el oxígeno
del alcohol ataca el azufre y su hidrógeno pasa a la base. No se rompe
ningún enlace con el carbono estereogénico, de modo que la disposición a
su alrededor es la misma. En la etapa SN2, el enlace C--O se rompe
mientras el \omterm{def:b1:electronic-effects:nucleophile}{nucleófilo} se une por el lado opuesto
(\cref{prop:b1:nucleophilic-substitution:sn2-inversion}).
\end{proof}

\begin{omfigure}
\begin{tikzpicture}[font=\small, >={Stealth}]
  \node (a) at (0,0) {\chemfig{C(-[:150]H_3C)(<[:210]H)(<:[:250]C_2H_5)-OH}};
  \node at (0,-1.4) {\cip{S}-butan-2-ol};
  \draw[->, thick] (1.3,0) -- (2.5,0) node[midway, above] {TsCl} node[midway, below, font=\scriptsize] {retención};
  \node (b) at (3.9,0) {\chemfig{C(-[:150]H_3C)(<[:210]H)(<:[:250]C_2H_5)-OTs}};
  \node at (3.9,-1.4) {tosilato \cip{S}};
  \draw[->, thick] (5.3,0) -- (6.6,0) node[midway, above] {\ce{CN-}} node[midway, below, font=\scriptsize] {SN2, inversión};
  \node (c) at (8.3,0) {\chemfig{[:0]NC-C(-[:30]CH_3)(<[:-30]H)(<:[:-70]C_2H_5)}};
  \node at (8.3,-1.4) {nitrilo \cip{R}};
\end{tikzpicture}

{\small Del \cip{S}-butan-2-ol al 2-metilbutanonitrilo: la tosilación
conserva la \omterm{def:b1:stereochemistry-in-depth:isomers}{configuración} y la etapa SN2 con cianuro la invierte. En
conjunto: OH sustituido por CN con inversión.}
\end{omfigure}

Al mismo objetivo puede llegarse en medio ácido: el grupo OH, protonado,
sale como agua, \ce{R-OH2+ -> R+ + H2O} (SN1, E1) o bajo un ataque (SN2).
Sin embargo, el ácido limita el \omterm{def:b1:electronic-effects:nucleophile}{nucleófilo} a los que lo resisten ---
iones haluro, agua, alcoholes --- y abre la puerta a los \omterm{def:b1:electronic-effects:intermediates}{carbocationes} y
sus reacciones secundarias.

\section{Conversión en halogenoalcanos}

\begin{proposition}[Los alcoholes y los haluros de hidrógeno]\label{prop:b1:alcohol-activation:hx-by-class}
Los alcoholes terciarios dan el halogenoalcano con ácido clorhídrico
concentrado a temperatura ambiente (SN1); los alcoholes primarios
necesitan ácido bromhídrico o yodhídrico y calentamiento (SN2 sobre el
alcohol protonado); los alcoholes secundarios reaccionan a velocidades
intermedias, a menudo por los dos mecanismos.
\end{proposition}

\begin{proof}
La protonación convierte el OH en \ce{-OH2+}. En un alcohol terciario, el
agua sale con facilidad y da un \omterm{def:b1:electronic-effects:intermediates}{carbocatión} estable que el haluro
captura (SN1): incluso el modesto \omterm{def:b1:electronic-effects:nucleophile}{nucleófilo} \ce{Cl-} basta. Un
\omterm{def:b1:electronic-effects:intermediates}{carbocatión} primario es demasiado inestable; el haluro debe empujar el
agua desde atrás (SN2), lo que exige un buen \omterm{def:b1:electronic-effects:nucleophile}{nucleófilo} (\ce{Br-},
\ce{I-}, mejores que \ce{Cl-} en medio prótico,
\cref{prop:b1:electronic-effects:nucleophilicity-basicity}) y calor.
\end{proof}

\begin{example}[Un ensayo de clasificación]\label{ex:b1:alcohol-activation:lucas}
Una mezcla de ácido clorhídrico concentrado y cloruro de cinc (un \omterm{def:b1:lewis-resonance-vsepr:lewis-acid}{ácido
de Lewis} que ayuda a salir al OH) enturbia un alcohol terciario al
instante --- el cloruro insoluble se separa ---, un alcohol secundario en
unos minutos y un alcohol primario apenas, a temperatura ambiente: el
orden de estabilidad de los \omterm{def:b1:electronic-effects:intermediates}{carbocationes}, hecho visible en un tubo de
ensayo.
\end{example}

\section{Deshidratación y formación de éteres en medio ácido}

\begin{definition}[Deshidratación de un alcohol]\label{def:b1:alcohol-activation:dehydration}
La \emph{deshidratación}\index{deshidratacion@deshidratación} de un alcohol es la
eliminación de agua, catalizada por un \omterm{def:b1:predominant-reaction:strong-acid}{ácido fuerte} (ácido sulfúrico o
fosfórico) y el calor, que da un alqueno:
$\mathrm{R_2CH{-}CR'_2OH} \to \mathrm{R_2C{=}CR'_2} + \ce{H2O}$.
\end{definition}

\begin{omfigure}
\setchemfig{atom sep=2.1em}
\schemestart
\chemfig{C(-[2]CH_3)(-[:210]H_3C)(-[:-30]CH_2CH_3)-OH}
\arrow{->[\ce{H+}]}
\chemfig{C(-[2]CH_3)(-[:210]H_3C)(-[:-30]CH_2CH_3)-OH_2^{+}}
\arrow{->[\ce{-H2O}]}
\chemfig{C^{+}(-[2]CH_3)(-[:210]H_3C)(-[:-30]CH_2CH_3)}
\schemestop

\medskip
\schemestart
\chemfig{C^{+}(-[2]CH_3)(-[:210]H_3C)(-[:-30]CH_2CH_3)}
\arrow{->[\ce{-H+}][(mayoritario)]}
\chemfig{C(-[2]CH_3)(-[:210]H_3C)=[:-30]CHCH_3}
\schemestop

\medskip
{\small \omterm{def:b1:alcohol-activation:dehydration}{Deshidratación} catalizada por ácido del 2-metilbutan-2-ol (E1):
protonación, pérdida de agua hasta el \omterm{def:b1:electronic-effects:intermediates}{carbocatión} terciario y pérdida de
un protón. Arrancar el protón del \ce{CH2} da el 2-metilbut-2-eno
trisustituido (Zaitsev, mayoritario); arrancarlo de un grupo metilo, el
2-metilbut-1-eno (minoritario).}
\end{omfigure}

\begin{proposition}[Éter o alqueno]\label{prop:b1:alcohol-activation:ether-vs-alkene}
Un alcohol primario calentado con ácido sulfúrico da sobre todo un éter
simétrico a temperatura moderada (una molécula sustituye a otra por
SN2), y sobre todo un alqueno a temperatura más alta; los alcoholes
secundarios y terciarios dan alquenos.
\end{proposition}

\begin{proof}
En un alcohol primario, el alcohol protonado puede ser atacado en el
carbono por una segunda molécula de alcohol (SN2, que da \ce{R-O-R} tras
la pérdida de un protón) o en un hidrógeno $\beta$ (E2). La eliminación
forma más moléculas de las que consume y se ve favorecida al subir la
temperatura (\cref{prop:b1:elimination:sn-vs-e}); la sustitución domina
cuando es más baja. Los alcoholes secundarios y terciarios se ionizan a
\omterm{def:b1:electronic-effects:intermediates}{carbocationes}, que pierden un protón con más facilidad de lo que un
\omterm{def:b1:electronic-effects:nucleophile}{nucleófilo} débil los captura, y tanto más al calentar; además, el agua
formada se elimina destilando el alqueno, volátil.
\end{proof}

\section{Ejercicios}

\begin{exercise}[$\star$]\label{exo:b1:alcohol-activation:1}
Escríbanse las reacciones del etanol con el sodio, con el hidruro de
sodio y con el hidróxido de sodio; calcúlese la constante de la última
($\mathrm{p}K_a$ del etanol, 15.9).
\end{exercise}

\begin{exercise}[$\star$]\label{exo:b1:alcohol-activation:2}
Dense los productos de: el metóxido de sodio con el 1-bromopropano; el
etóxido de sodio con el yodometano; el fenóxido de sodio con el
bromoetano.
\end{exercise}

\begin{exercise}[$\star$]\label{exo:b1:alcohol-activation:3}
Escríbanse la tosilación del propan-1-ol y la reacción del \omterm{def:b1:alcohol-activation:sulfonate}{tosilato} con
yoduro de sodio. ¿Qué ventaja tiene sobre la reacción directa del alcohol
con yoduro de sodio?
\end{exercise}

\begin{exercise}[$\star$]\label{exo:b1:alcohol-activation:4}
Dense los alquenos que se forman por \omterm{def:b1:alcohol-activation:dehydration}{deshidratación} ácida del butan-2-ol y
del 2-metilpentan-2-ol, y el mayoritario.
\end{exercise}

\begin{exercise}[$\star\star$]\label{exo:b1:alcohol-activation:5}
Propóngase una \omterm{def:b1:alcohol-activation:williamson}{síntesis de Williamson} del 1-etoxi-2-metilpropano y
explíquese la elección de los participantes. ¿Por qué la otra elección
daría un alqueno?
\end{exercise}

\begin{exercise}[$\star\star$]\label{exo:b1:alcohol-activation:6}
El \cip{R}-octan-2-ol se convierte en su \omterm{def:b1:alcohol-activation:sulfonate}{tosilato} y después se trata con
etóxido de sodio en etanol. Dese la \omterm{def:b1:stereochemistry-in-depth:isomers}{configuración} del éter. ¿Qué se
obtendría calentando el mismo alcohol con ácido sulfúrico?
\end{exercise}

\begin{exercise}[$\star\star$]\label{exo:b1:alcohol-activation:7}
Escríbase el mecanismo de la reacción del 2-metilpropan-2-ol con ácido
clorhídrico concentrado y explíquese por qué la reacción es rápida a
temperatura ambiente.
\end{exercise}

\begin{exercise}[$\star\star$]\label{exo:b1:alcohol-activation:8}
Escríbanse el mecanismo de la formación del etoxietano a partir del
etanol en ácido sulfúrico y la eliminación que compite con ella.
\end{exercise}

\begin{exercise}[$\star\star$]\label{exo:b1:alcohol-activation:9}
En el ensayo de clasificación de este capítulo, ¿qué alcohol se enturbia
primero entre el butan-1-ol, el butan-2-ol y el 2-metilpropan-2-ol?
Escríbase el producto.
\end{exercise}

\begin{exercise}[$\star\star\star$]\label{exo:b1:alcohol-activation:10}
El 3,3-dimetilbutan-2-ol, calentado en medio ácido, da sobre todo
2,3-dimetilbut-2-eno, cuyo esqueleto carbonado difiere del del alcohol.
Propóngase un mecanismo con una migración 1,2 de un grupo metilo en el
\omterm{def:b1:electronic-effects:intermediates}{carbocatión}, y justifíquese la migración.
\end{exercise}

\begin{exercise}[$\star\star\star$]\label{exo:b1:alcohol-activation:11}
Diséñese una síntesis del \cip{S}-2-metoxibutano a partir del
\cip{R}-butan-2-ol, y otra a partir del \cip{S}-butan-2-ol, controlando en
cada una la \omterm{def:b1:stereochemistry-in-depth:isomers}{configuración}.
\end{exercise}

\begin{exercise}[$\star\star\star$]\label{exo:b1:alcohol-activation:12}
Demuéstrese, con el \omterm{def:b1:extent-q-and-k:quotient}{cociente de reacción} (\cref{ch:b1:extent-q-and-k}),
que el equilibrio \ce{2CH3CH2OH <=> (CH3CH2)2O + H2O} puede desplazarse
hacia el éter eliminando continuamente el agua.
\end{exercise}

\section{Problema: dos caminos hacia un éter}

\begin{problem}[{Problema de fin de semana --- las dos vías de Williamson
hacia el metil terc-butil éter, un éter con la estereoquímica controlada a
partir de un alcohol ópticamente activo, la deshidratación de un alcohol
terciario y la masa de éter obtenida}]
\label{pb:b1:alcohol-activation:1}
El 2-metoxi-2-metilpropano (metil terc-butil éter, MTBE) se producía a
gran escala como aditivo de combustibles. Masas molares (\unit{g/mol}):
\ce{C4H10O}
74.0, \ce{C5H12O} 88.0, \ce{CH3I} 141.9.

\medskip
\noindent\textbf{Parte I --- El MTBE por la \omterm{def:b1:alcohol-activation:williamson}{síntesis de Williamson}.}
\begin{enumerate}
  \item Escríbanse las dos parejas de participantes (\omterm{def:b1:alcohol-activation:alkoxide}{alcóxido} y haluro)
        que podrían dar MTBE.
  \item ¿Qué pareja falla? Escríbase la reacción que tiene lugar en su
        lugar.
  \item Escríbanse la reacción que funciona y su mecanismo.
  \item ¿Cómo se prepara el terc-butóxido de sodio a partir del
        2-metilpropan-2-ol? ¿Por qué no con hidróxido de sodio?
  \item ¿Por qué se lleva a cabo la reacción en el alcohol seco o en un
        \omterm{def:b1:intermolecular-forces-solvents:solvent-classes}{disolvente aprótico}?
  \item ¿Es \omterm{def:b1:stereochemistry-in-depth:stereocentre}{quiral} el producto?
\end{enumerate}

\medskip
\noindent\textbf{Parte II --- Un éter ópticamente activo.}
\begin{enumerate}[resume]
  \item El \cip{R}-butan-2-ol se convierte en su \omterm{def:b1:alcohol-activation:sulfonate}{tosilato}. Escríbase la
        reacción y dese la \omterm{def:b1:stereochemistry-in-depth:isomers}{configuración} del \omterm{def:b1:alcohol-activation:sulfonate}{tosilato}.
  \item El \omterm{def:b1:alcohol-activation:sulfonate}{tosilato} reacciona con metóxido de sodio. Escríbanse el
        producto y su \omterm{def:b1:stereochemistry-in-depth:isomers}{configuración}.
  \item ¿Qué \omterm{def:b1:stereochemistry-in-depth:isomers}{configuración} del alcohol da el \cip{R}-2-metoxibutano por la
        misma vía?
  \item Otra vía: el \omterm{def:b1:alcohol-activation:alkoxide}{alcóxido} del \cip{R}-butan-2-ol con yodometano.
        ¿\omterm{def:b1:stereochemistry-in-depth:isomers}{Configuración} del producto? ¿Por qué?
  \item ¿Qué reacción secundaria compite en la pregunta 8, y por qué está
        limitada aquí?
  \item ¿Por qué calentar el \cip{R}-butan-2-ol con metanol en ácido
        sulfúrico daría un éter casi racémico, si es que lo da?
\end{enumerate}

\medskip
\noindent\textbf{Parte III --- \omterm{def:b1:alcohol-activation:dehydration}{Deshidratación} de un alcohol terciario.}
\begin{enumerate}[resume]
  \item Escríbase el mecanismo de la \omterm{def:b1:alcohol-activation:dehydration}{deshidratación} ácida del
        2-metilbutan-2-ol.
  \item Dense los dos alquenos y predígase el mayoritario.
  \item ¿Por qué no se forma ningún éter a partir de este alcohol?
  \item ¿Por qué se destila el 2-metilbut-2-eno a medida que se forma?
  \item Dibújese el \omterm{def:b1:elementary-steps:energy-profile}{perfil de energía} de la etapa E1 y márquese la \omterm{def:b1:elementary-steps:rds}{etapa
        determinante de la velocidad}.
  \item ¿Podría deshidratarse un alcohol primario por el mismo mecanismo?
\end{enumerate}

\medskip
\noindent\textbf{Parte IV --- Masas.}
\begin{enumerate}[resume]
  \item Se convierten \qty{10.0}{g} de 2-metilpropan-2-ol en el \omterm{def:b1:alcohol-activation:alkoxide}{alcóxido}.
        Calcúlese la cantidad.
  \item ¿Qué masa de yodometano hace falta para un exceso del
        \qty{10}{\percent}?
  \item Calcúlese la masa teórica de MTBE.
  \item Calcúlese la masa obtenida con un \omterm{def:b1:extent-q-and-k:fractional-extent}{rendimiento} del
        \qty{75}{\percent}.
\end{enumerate}
\end{problem}
